\section{Thin triangles and escape windows}\label{sec:skinny}

At the thin end of the family, the small mass $B=\epsilon$ is far from a
tight pair of masses $A=\sqrt{1-\epsilon^2}$ and 1, initially separated by
$\epsilon$. Two clocks then coexist: the inner pair moves on a time scale
$\epsilon^{3/2}$, while the outer body returns on a time scale of order
one. A single outer excursion contains an unbounded number of binary
cycles as $\epsilon\to0$. Exclusion of one early encounter therefore
leaves an increasing number of later passages to control.

\subsection{The first close encounter}

For $X=q_3-q_1$, put $X=\epsilon x$ and
$t=\epsilon^{3/2}\tau/\sqrt{A+1}$, and rotate so that $x(0)=1$.
The leading motion is radial Kepler fall. In Levi--Civita coordinates
$x=z^2$, $d\tau=|z|^2ds$, its solution is
$z_0(s)=\cos(s/\sqrt2)$, with collision at $s_c=\pi/\sqrt2$.
The true parameter-dependent equations are analytic at this comparison
collision. Their first transverse variation appears only at order five:
\[
 v_{ss}+\tfrac12v=-\tfrac38\cos^5(s/\sqrt2),\qquad
 v(0)=v_s(0)=0,\qquad v(s_c)=-{15\pi\over128}.
\]
Consequently the physical first encounter misses collision by
\begin{equation}\label{eq:skinny-miss}
 r_{13,\min}={225\pi^2\over16384}\epsilon^{11}(1+O(\epsilon)).
\end{equation}
The proof uses the simple zero of $\operatorname{Re}z_0$ and the nonzero
fifth-order imaginary displacement; away from that zero compact continuity
keeps all pairs separated. Squaring the regularized displacement and
restoring the factor $\epsilon$ gives the eleventh power. The symbolic
coefficient is also obtained directly from the oscillator Green function:
with $a=s/\sqrt2$,
\[
 v(s_c)=-{3\over4}\int_0^{\pi/2}\cos^6 a\,da
 =-{15\pi\over128}.
\]
The same calculation gives the specific pair angular momentum
$\ell_{13}=-(15\pi/64)\epsilon^{11/2}+O(\epsilon^{13/2})$.

The first passage is not an escape certificate. Its outer radial escape
margin is $-2+O(\epsilon)$, while the pair remains tightly bound. The
subsequent return must be analyzed.

\subsection{The restricted problem at the outer plunge}\label{sec:restricted}

Let $X=q_3-q_1$ and $Y=q_2-(Aq_1+q_3)/(A+1)$. Near the outer
plunge, set $X=\epsilon R$, $Y=\epsilon Z$, and
$t=t_*+\epsilon^{3/2}\theta$. With $\Phi(v)=v/|v|^3$, the limiting
equations are
\begin{equation}\label{eq:restricted}
 R_{\theta\theta}=-2\Phi(R),\qquad
 Z_{\theta\theta}=-\Phi(Z+R/2)-\Phi(Z-R/2).
\end{equation}
The two heavy masses have become equal, and the light particle no longer
affects their motion. On the invariant symmetric subspace,
$R=(r,0)$, $Z=(0,z)$, and
\[
 r=\cos^2\psi,\qquad d\theta=r\,d\psi,\qquad
 z_{\theta\theta}=-{2z\over(z^2+r^2/4)^{3/2}}.
\]
Here the binary has apocenter one and is continued through its collisions
in Levi--Civita coordinates. All limiting arguments through these
collisions require the light particle to remain separated.

Release the light particle from $z=0$ at binary apocenter with speed $v$.
Small positive speeds turn, whereas sufficiently large speeds escape.
The boundary of the turning set contains a parabolic orbit, tending to
infinity with zero speed; reflection and time reversal extend it to a
parabolic-to-parabolic orbit. For $z>0$, $z_{\theta\theta}<0$ prevents a
bounded nonturning boundary orbit. Openness of finite turning and
positive-speed escape excludes both alternatives at the boundary.
The speed $v_*$ satisfies $\sqrt8<v_*<4$; interval integration with the
analytic tail bound
\[
 0\le H_K(\infty)-H_K(z)\le{1\over4z^3},\qquad
 H_K={z_\theta^2\over2}-{2\over z},
\]
refines this to $2.9051113<v_*<2.9051116$.

\subsection{Matching the binary clock and the full incoming state}

Write $M=A+1$, $N=M+\epsilon$, and
$\rho_0^2=2A^2/M$. Replacing the inner pair by its combined point mass
provides an explicit reference phase:
\begin{equation}\label{eq:clock}
 \Phi_{\rm ref}(\epsilon)
 ={\pi\over\epsilon^{3/2}}\sqrt{M\over N}\,\rho_0^{3/2}
 =\pi\epsilon^{-3/2}-{\pi\over4}\epsilon^{-1/2}
   -{15\pi\over32}\epsilon^{1/2}+O(\epsilon^{3/2}).
\end{equation}
The reference clock approximates the asymptotic phase intercept of the
interacting system. To transfer an escape window, one needs convergence
of the full incoming state as well as its phase.

\begin{theorem}[Incoming-state matching]
On a fixed sufficiently large scaled incoming section $\Sigma_K^-$, let
$\Gamma_K^-(\chi)$ denote the incoming parabolic state of the equal-heavy-mass
restricted problem, parameterized by its asymptotic mean-anomaly intercept
$\chi$. As $\epsilon\to0$, either a classical collision ends the exact tied
orbit before it reaches the section, or its full regularized section state
$S_\epsilon(K)$ satisfies
\[
 \operatorname{dist}\bigl(S_\epsilon(K),
 \Gamma_K^-(\Phi_{\rm ref}(\epsilon))\bigr)=o(1).
\]
\end{theorem}
\begin{proof}
Stop the outer infall at physical radius $\rho=\epsilon^\alpha$ with
$0<\alpha<1/6$. Cancellation of the outer dipole gives a monopole-force
remainder $O(\epsilon^2\rho^{-4})$. Inner-energy control, variation of the
regularized oscillator, and comparison of its physical clock then give
an incoming intercept error $O(\epsilon^{1/2-3\alpha})=o(1)$.
On the remaining tail, write $y=|Z|$, $y_e=\epsilon^{\alpha-1}$,
$\mathcal H_N=|Z_\theta|^2/2-N/y$, and
$\mathcal L=Z\times Z_\theta$. The dipole cancellation and a simultaneous
energy and angular-momentum bootstrap give, uniformly for $K\le y\le y_e$,
\[
 |\mathcal H_N-E_{\rm ref,\epsilon}|\le Cy^{-3},\qquad
 |\mathcal L|\le Cy^{-3/2},\qquad -y_\theta\asymp y^{-1/2},
\]
where $E_{\rm ref,\epsilon}=-N\epsilon/\rho_0$ is the scaled monopole
energy. Binary control through every close passage uses the oscillator
invariants
\[
 \mathcal J=(E+1,L_b,K_b),\quad
 L_b=z_xp_y-z_yp_x,\quad K_b=p_xp_y-(E/2)z_xz_y,
\]
with $p=z_s$. On the radial Kepler circle the normal Jacobian
$\det\partial(L_b,K_b)/\partial(z_y,p_y)=-1/2$ stays nonzero, including
at binary collision. The forced equations give
$|\mathcal J_s|\le C\epsilon y^{-3}$; summing complete LC blocks and
bounding the two incomplete end blocks gives a clock error
$O(\epsilon^{1-3\alpha/2})+O(\epsilon|\log\epsilon|)=o(1)$.

Retain the finite-$\epsilon$ monopole flight-time subtraction until $y$
is fixed. Its phase error is at most $Cy^{-1/2}+o(1)$, which includes
the finite quadrupole correction from infinity. Compactness in LC
variables gives a limiting restricted tail. Its transverse displacement
has the inherited bound $O(y^{-1})$, whereas the homogeneous transverse
modes grow as $y$ and $y^{1/2}$. Thus the limit is rectilinear. Taking
$\epsilon\to0$ at fixed $y$ and then $y\to\infty$ identifies the
intercept as $\Phi_{\rm ref}(\epsilon)$ in the limiting sense.
Uniqueness of the incoming parabolic graph gives the stated full-state
convergence. The uniform tail estimates and their LC block bounds are
recorded in \path{docs/INCOMING_TAIL.md}.
\end{proof}

\subsection{Transversality and the escape arcs}

At $z=0$, let $V_u(\phi)$ and $V_s(\phi)$ be the positive crossing
speeds on the incoming and outgoing parabolic curves, with $\phi=0$ at
binary apocenter. Reversibility gives $V_s(\phi)=V_u(-\phi)$.
The transversality condition is $V_u'(0)\ne0$. Fixing the phase is
essential: the time-translation Jacobi field decays at parabolic infinity
and is not removed merely by prescribing zero limiting data.

For the longitudinal Jacobi field $h$, put $q=h/v$, $w=z_\theta$,
$p=q_\theta$, and $d^2=z^2+r^2/4$. The finite regular system is
\[
 z_\psi=rw,\quad w_\psi=-{2rz\over d^3},\quad
 q_\psi=rp,\quad p_\psi=r{4z^2-r^2/2\over d^5}q,
 \qquad(z,w,q,p)(0)=(0,v,1,0).
\]
An order-20 interval Taylor calculation on 256 overlapping speed slabs
proves
\[
 p(\pi/2)>1/125\qquad\hbox{uniformly for }14/5\le v\le4.
\]
Here $v$ denotes launch speed, not the near-isosceles parameter.
Force comparison gives $h>0$ through this first contraction and a positive
Jacobi coefficient thereafter. Thus the certified derivative forces
$h(\theta)$ to grow at least linearly.

To relate this growth to transversality, compactify infinity by $z=2/x^2$.
The exact field is analytic at $x=0$, and its time-$\pi$ map in binary
eccentric anomaly, after $X=x$, $Y=w-\sqrt2x$, is
\[
 X_1=X-{\pi\over8}X^3(Y+\sqrt2X)+O_7,\qquad
 Y_1=Y+{\pi\sqrt2\over8}X^3Y+O_7.
\]
The degree-four terms contract the positive $X$ axis and expand its normal
direction; $O_7$ denotes analytic terms of total degree at least seven.
McGehee's degenerate stable-manifold theorem gives a parabolic
graph $Y=\varphi(X)$ with $\varphi'(X)\to0$~\cite{McGehee1973}.
Along it, $X_n^{-3}\asymp n$ and every tangent tends to zero.
If the parabolic curves were tangent at $\phi=0$, reversibility would
make both slopes vanish. The difference of the time-shift tangent and
the parabolic tangent with fixed phase would then be the field $h$ above.
But $\delta X_n=-X_n^3h(\theta_n)/4$ would stay bounded away from zero,
contradicting tangent contraction. This proves transversality.

Consequently $V_u-V_s$ changes sign linearly at zero. The parabolic graph
separates an open hyperbolic-escape side from an open finite-turn side.
This is the source of the escape-intercept arcs used below.

\subsection{Transfer to integer triangles}

Fix a compact subarc $J$ of positive length strictly inside a restricted escape window.
Incoming-state matching and continuous dependence carry sufficiently
small exact tied members with $\Phi_{\rm ref}(\epsilon)\bmod2\pi\in J$
to a finite outgoing section. There the binary specific energy and outer
speed have the forms
\[
 e=-C/\epsilon+o(\epsilon^{-1}),\qquad
 \dot\rho=s\epsilon^{-1/2}+o(\epsilon^{-1/2}),\qquad C,s>0,
\]
with strict margins uniform on the compact arc. At a sufficiently distant
scaled section, Lemma~\ref{lem:escape} applies. An earlier classical
collision already excludes periodicity; every other member escapes or
subsequently collides. This proves all-time exclusion on these windows,
not just on the first excursion.

\begin{theorem}[Infinitely many thin-triangle intervals]
There are infinitely many open intervals of real $u$ accumulating at zero
on which the exact tied solution is nonperiodic. Each contains rational
parameters and therefore primitive integer Pythagorean triples.
\end{theorem}
\begin{proof}
Equation~\eqref{eq:clock} gives
$\Phi_{\rm ref}'(\epsilon)\sim-(3\pi/2)\epsilon^{-5/2}$.
The preimages of the interior of $J$ therefore include infinitely many
disjoint intervals accumulating at zero. The uniform transfer above applies
to every sufficiently small one. Finally $B(u)$ is a homeomorphism near
zero and rational $u$ are dense. Different parameters in the fundamental
interval give different primitive triples after scaling.
\end{proof}

\begin{theorem}[A positive-density explicit family]
For a set of positive lower natural density of integers $n\ge1$, the
primitive triple
\[
 (a_n,b_n,c_n)=(4n^2-1,4n,4n^2+1)
\]
has a nonperiodic Pythagorean--Burrau solution.
\end{theorem}
\begin{proof}
Here $u_n=1/(2n)$ and $B_n=4n/(4n^2+1)$. The exact reference clock gives
\[
 F(x):={\Phi_{\rm ref}(B_x)\over2\pi}
 ={(2x-1)^{3/2}(2x+1)(4x^2+1)^{3/4}\over32x^{5/2}},
\]
with
\[
 F''(x)={3\over8}x^{-1/2}+{1\over32}x^{-3/2}+O(x^{-5/2}).
\]
The differentiated remainder follows from analyticity in $1/x$ after
factoring out $x^{3/2}$. For each nonzero integer $h$, van der Corput's
second-derivative estimate bounds the Weyl sum on $N<n\le2N$ by
$O_h(N^{3/4})$. Dyadic summation and Weyl's criterion give equidistribution
of $F(n)$ modulo one. Thus the compact arc $J$ selects indices of density
$|J|/(2\pi)>0$, and the escape transfer proves nonperiodicity for every
sufficiently large selected index. Since $\gcd(4n,4n^2-1)=1$, the triples
are primitive. This density concerns the displayed family of indices,
not the set of all primitive Pythagorean triples ordered by hypotenuse.
\end{proof}

The returning side requires additional control through successive central
encounters. The next section gives the available first-turn exclusion,
the near-triple-collision reduction, and the local collision classifications.
